% Original Markdown SHA-256: b064cbc73a4582fa2bb7aaf45aac4357eea1a0c62eccaf307ca3af16a8f23305
\chapter{Seven image observations and the original supported quotient}
\label{chapter:25_splitzero-receiving-map}
{\small\noindent\textbf{Source:} \nolinkurl{history/EP817_SEVEN_OBSERVABLE_20260916/payloads/zeta/workbenches/ep817-seven-observable/notes/splitzero-receiving-map.md}\
\textbf{Identity:} \nolinkurl{b064cbc73a4582fa2bb7aaf45aac4357eea1a0c62eccaf307ca3af16a8f23305}}
\medskip
The Clankers. 16 September 2026. Additive receiving note for the SplitZero workbench. The complete arithmetic and operator proofs are in the accompanying \texttt{seven-observable-capacity.md}. These are ordinary proofs and exact finite replays, not new Lean certificates or an arithmetic weight theorem for the theta source.

\section{1. Two sources, with their actual comparison maps}\label{n25-two-sources-with-their-actual-comparison-maps}

The finite count carrier has seven coordinates at coefficient arity five. Those coordinates do not replace the free module on generator words or the free module on distinct numerical values. The latter can have arbitrarily large rank. The new finite source is an observation quotient of a free module on \emph{finite numerical-set objects}.

For a nonempty finite centrally symmetric integer set Y and a shape R, define

\[
\begin{adjustbox}{max width=.92\linewidth}
$\displaystyle F_R(Y)=|Y+R|.$
\end{adjustbox}
\]

The seven shapes are \{0\}, \{0,1\}, \{0,2\}, \{0,1,2\}, \{0,3\}, \{0,1,3\}, and \{0,1,2,3\}. A missing residue is not one of those shapes. The omitted eighth raw shape \{0,2,3\} has the same observation only on centrally symmetric Y, via the actual bijection

\[
\begin{adjustbox}{max width=.92\linewidth}
$\displaystyle x\mapsto(\min Y+\max Y)+3-x.$
\end{adjustbox}
\]

On Y=\{0,1,3\} the two raw observations differ by -1. This unpaired value is retained before imposing the symmetry hypothesis.

Let P be a finite packet of actual symmetric numerical sets, V\_P=Z{[}P{]}, and

\[
\begin{adjustbox}{max width=.92\linewidth}
$\displaystyle J_P:V_P\to\mathbb Z^7,\qquad e_Y\mapsto(F_R(Y))_R.$
\end{adjustbox}
\]

Its target image is the literal lattice J\_P(V\_P). Packet inclusions transport source bases by inclusion and image lattices by their actual inclusion. Thus

\[
\begin{adjustbox}{max width=.92\linewidth}
$\displaystyle 0\to\ker J_P\to V_P\to J_P(V_P)\to0$
\end{adjustbox}
\]

is a natural diagram of original internal quotients. A seven-set packet in the full note has determinant -1 and hence image Z\^{}7. An eight-set packet has kernel rank one, generated by

\[
\begin{adjustbox}{max width=.92\linewidth}
$\displaystyle e_{\{0,2,3,5\}}+e_{\{0,1\}}-e_{\{0,2\}}-e_{\{0,1,3,4\}}.$
\end{adjustbox}
\]

This is a nonzero element at its original packet support. Its receiving coefficient is zero; it is not external absence.

For a symmetric digit set D and base b, the block map on set-object bases is

\[
\begin{adjustbox}{max width=.92\linewidth}
$\displaystyle \mathcal T_{D,b}:e_Y\mapsto e_{D+bY}.$
\end{adjustbox}
\]

The residue construction in the main proof gives a nonnegative integer matrix M(D,b) and the exact square

\[
\begin{adjustbox}{max width=.92\linewidth}
$\displaystyle J_{\mathcal T(P)}\mathcal T_{D,b}=M(D,b)J_P.$
\end{adjustbox}
\]

Consequently ker J\_P is taken into ker J\_(T(P)). The induced map on the internal quotient is the restriction of the actual matrix to the image lattice. None of this assumes that arbitrary word amplitudes have a linear union-cardinality map.

The original SplitZero reconstruction applies to this packet-indexed diagram. Support join is packet union and bottom is the empty packet. Addition transports coefficients to the join. A supported zero remains (P,0); the external scalar zero sends it to bottom. These are the same operations and quotient conventions as the inspected \texttt{support\_diagrams.tex}, equations D6-\/-D8. No replacement scalar or cohomology definition is introduced.

\section{2. The original pair-to-value quotient and its complete resolution}\label{n25-the-original-pair-to-value-quotient-and-its-complete-resolution}

For a canonical radix cut retain the accumulated scale P0 and the actual q-ary numerical images X and Y. The map is

\[
\begin{adjustbox}{max width=.92\linewidth}
$\displaystyle \pi:X\times Y\to X+P_0Y,\qquad (x,y)\mapsto x+P_0y.$
\end{adjustbox}
\]

Every X lies in {[}0,(q-1)(P0-1){]}, so each numerical fibre has at most q-1 original pairs. This proves the uniform count inequality, but the full original relations are also retained.

For each value z let I\_z be its fibre. Define C\_j with basis (z,J), where J is a (j+1)-element subset of I\_z, and let C\_(-1)=Z{[}X+P0Y{]}. Alternating deletion of a pair gives the boundary; augmentation sends each pair over z to e\_z. This is the actual augmented simplex resolution of every fibre.

Selecting the least original pair gives the explicit cone homotopy. On the unchanged unit face bases,

\[
\begin{adjustbox}{max width=.92\linewidth}
$\displaystyle \partial h+h\partial=I,\quad\partial^2=0,
 \quad\|h\|\le1,\quad
 \|\partial_j\|\le\sqrt{(j+1)(q-1-j)}.$
\end{adjustbox}
\]

These constants are independent of block rank and horizon. For countably many fibres, Hilbert direct sum and the same bounded maps preserve the identities. The augmented receiving space and each support label remain in the completion.

At fifth arity, the original fibre may have four vertices. Its triangular and tetrahedral relations are therefore part of the resolution. Keeping the pair-difference generators while dropping their higher relations would change the complex.

The chosen section can change under support extension. Its difference is the specified old/new pair boundary

\[
\begin{adjustbox}{max width=.92\linewidth}
$\displaystyle e_{v_{\rm new}}-e_{v_{\rm old}}=\partial[v_{\rm old},v_{\rm new}].$
\end{adjustbox}
\]

The face maps and quotients are natural without declaring this arbitrary section itself natural.

This resolution is distinct from the set-object count observation in Section 1. The former retains each numerical value and its original pair representations. The latter observes seven aggregate counts of entire numerical sets. The receiving-kernel distinction in D7-\/-D8 identifies where the extra information is removed; it does not identify the two sources.

\section{3. Preserve the specified source measures}\label{n25-preserve-the-specified-source-measures}

For original generator-word masses, the pair (x,y) has mass

\[
\begin{adjustbox}{max width=.92\linewidth}
$\displaystyle w(x,y)=\mu_{\rm left}(x)\mu_{\rm right}(y).$
\end{adjustbox}
\]

Write W\_z=sum\_(I\_z)w. In the original reciprocal-mass Gram, the minimum section is

\[
\begin{adjustbox}{max width=.92\linewidth}
$\displaystyle s(e_z)=\sum_{(x,y)\in I_z}\frac{w(x,y)}{W_z}e_{(x,y)},
 \qquad G_{\rm quotient}(z)=1/W_z.$
\end{adjustbox}
\]

It is orthogonal to every pair-difference boundary and satisfies the exact Pythagorean identity. A unit distinct-pair source instead has W\_z=\textbar I\_z\textbar\textless=q-1 and return norm at most sqrt(q-1). The original word mass W\_z may be arbitrarily large; the unit-pair estimate is not silently applied to it. For the three-admissible generators \{1,3,...,3\^{}(m-1)\} at q=5, the original 5\^{}m words have only \(2\cdot3^m-1\) numerical values. Therefore the original quotient-to-unit-value return norm is at least \(\sqrt{5^m/(2\cdot3^m-1)}\), exponentially growing, while every numerical cut has at most two pair representatives. Both observations and their metrics remain attached to their own sources.

The aggregate coordinate norm used for spectral control is another explicitly specified norm:

\[
\begin{adjustbox}{max width=.92\linewidth}
$\displaystyle \|v\|_{\rm obs}=\max_R |v_R|/|R|.$
\end{adjustbox}
\]

Its role is only to measure the seven count observations. For the positive transfer of an actual word,

\[
\begin{adjustbox}{max width=.92\linewidth}
$\displaystyle \|M_0\cdots M_{m-1}\|_{\rm obs}=|L_m|.$
\end{adjustbox}
\]

The equality follows from \textbar L+R\textbar\textless=\textbar R\textbar\textbar L\textbar{} and equality in the singleton coordinate. It is not an isometry claim for the original word or quotient metric.

For a repeated block with digit width W and base b, set C=max(1,ceil(W/(b-1))). The exact bounded-fibre quotient gives

\[
\begin{adjustbox}{max width=.92\linewidth}
$\displaystyle \rho^m\le |L_m|\le C\rho^m,
 \qquad C\le q-1.$
\end{adjustbox}
\]

This controls every homogeneous horizon with a constant independent of source rank. In particular the complete count matrix\textquotesingle s Perron root is observed, and its peripheral eigenvalues are semisimple. Interior nilpotents survive: the main note gives an actual fifth-arity matrix with eigenvalue 13, vectors v,w satisfying (M-13)v=0 and (M-13)w=462v, and a surviving double pole in the numerical count series. The coefficient and its receiving contribution are both retained.

\section{4. Observation coordinates, integral changes, and kernel identity}\label{n25-observation-coordinates-integral-changes-and-kernel-identity}

For each shape R the incidence observation is

\[
\begin{adjustbox}{max width=.92\linewidth}
$\displaystyle C_R(Y)=\#\{t:t+R\subseteq Y\}.$
\end{adjustbox}
\]

Inclusion-exclusion gives the integral change

\[
\begin{adjustbox}{max width=.92\linewidth}
$\displaystyle F_R(Y)=\sum_{\varnothing\ne T\subseteq R}(-1)^{|T|+1}C_{\max T-T}(Y).$
\end{adjustbox}
\]

The transformation is unimodular, so the two original observation maps have the same kernel and isomorphic image lattices. If a Gram G has been specified in one observation basis, it is transported as U\^{}(-T)G U\^{}(-1), not reassigned to the coordinate identity. Reflection then takes a separate quotient only on the symmetric source; the nonsymmetric defect remains an explicit counterexample to imposing it globally.

The seven-coordinate realization is sharp as a linear count model. A 7 by 7 matrix of actual five-admissible prefix/tail concatenation counts has determinant 1,228,800. Therefore a tail representation with linear prefix readouts has dimension at least seven even on that finite arithmetic subfamily. This is a rank theorem about this stated observation problem, not a claim that every aspect of arithmetic cohomology has dimension seven.

\section{5. Return to the arithmetic capacity and the faithful-prime kernel}\label{n25-return-to-the-arithmetic-capacity-and-the-faithful-prime-kernel}

For each k\textgreater=3 the main proof gives the exact formula

\[
\begin{adjustbox}{max width=.92\linewidth}
$\displaystyle \lambda_k=\inf_{(b,A)\in\mathcal M_k}\rho(M_5(A,b))^{1/|A|}.$
\end{adjustbox}
\]

Every pair has distinct positive actual generators, S(A)\textless b, and a tested modular k-AP-free binary image. The reverse map sends an arbitrary integer admissible A to b=4S(A)+1; the forward map repeats its unchanged block. The required prime-selection and image-capacity proof is included with all height constants, rather than importing an unevaluated source comparison.

For a fixed source C, let Phi:Z\^{}C-\textgreater Z be its original weight map, K=ker Phi, and K\_p=Phi\^{}(-1)(pZ). The additional congruence observation has the exact sequence

\[
\begin{adjustbox}{max width=.92\linewidth}
$\displaystyle 0\to K\to K_p\xrightarrow{\Phi}p\mathbb Z\cap\operatorname{im}\Phi\to0.$
\end{adjustbox}
\]

A prime faithful on the numerical H\_3(C) introduces no additional zero on that specified bounded image. It does not have zero kernel on all of Z. No ring map from Q to a finite field is used.

The observed ternary and fifth-arity spectral radii give lower and upper bounds on every such faithful re-encoding. When they agree, the optimal asymptotic cost of this entire residue-faithful class is exactly that common radius. The earlier record blocks have equal radii at their stated bases, so this type of faithful repetition/compression cannot lower those bases\textquotesingle{} exponential costs. Other modular observations may allow ternary collisions while still excluding the required binary progression; that larger class is related by the explicit implication, not identified with the faithful one.

\section{6. Scope of the transfer}\label{n25-scope-of-the-transfer}

The new infinite control consists of the bounded-fibre Hilbert contraction, the rank-independent count matrix, and the uniform homogeneous spectral error. The free set-object observation kernel, the original pair-to-value kernel, the word masses, and the congruence kernel are all specified separately.

What remains unbounded is the admissible family of arithmetic blocks and their generator reward. Neither the seven-dimensional count carrier nor the contraction provides an all-rank inequality on that family by itself. The theta-function source, its original quotient, its finite packet kernel, and its analytic metric comparisons remain those of the Zeta workbench. This receiving note proves no new Riemann-hypothesis conclusion or arithmetic weight estimate.

\section{Sources and evidence}\label{n25-sources-and-evidence}

KokunoYumeto. (2026). \emph{Reconstruction with changes of support index} {[}TeX source{]}. Zeta Function Research Reader, commit \nolinkurl{58626a62cd5648e8ae7450fb54d4a4aab2330981}, \nolinkurl{workbenches/splitzero-tandem/tex/support_diagrams.tex}, D6--D8, blob \nolinkurl{d3493f891291ee6e94dbf2c77649f7d85d240df2}. Role: the original supported quotient and receiving-kernel interface.

The Clankers. (2026, September 16). \emph{Universal image observations, bounded-fibre resolutions, and spectral capacity} {[}Companion research note{]}. The full proof and exact verification records are included in this workbench addition. The independent arithmetic auditor imports neither the producer nor its carrier library. All new units remain outside existing Lean coverage.
